\scanpage{145}
déduit de la structure d'extension. \deleted{Cet homoph.} (5) est compatible avec les op. du groupe $\mathbb Z/2\mathbb Z$, et définit donc un iso
\[
\underline{\mathrm{Arf}}'\simeq\underline{\mathrm{Arf}}\mathbin{\wedge}_{(\mathbb Z/2)_S}\mu_{2S}\tag*{(7)}
\]
--- donc, si $2$ inv. sur $S$, un iso.
\[
\underline{\mathrm{Arf}}\xrightarrow{\sim}\underline{\mathrm{Arf}}'=\underline M((-1)^n\delta(q^{*})^{-1})
\]

\deleted{f) Les $\underline{G}_H$ sont des espaces homogènes sous $\mathrm O(q)$}

\begin{center}\rule{.24\linewidth}{.4pt}\\[-.5ex]\rule{.15\linewidth}{.4pt}\end{center}

\uline{Application aux groupes orthogonaux}

\uline{Thm} a) Supposons $N=2n+1$, Alors
\[
\underline{\mathrm O}(q)\simeq\mu_2\times\underline{\mathrm{SO}}(q)
\]
et $\underline{\mathrm{SO}}(q)$ est \inserted{(affine)} lisse à fibres conn.

\deleted{b) $\underline{\mathrm{SO}}(q)$ \unsure{b)} $\mathrm{O}(q)$} opère sur chaque $\underline{G}_H(E,q)$ \inserted{via $\underline{\mathrm{SO}}(q)\simeq\underline{\mathrm O}(q)/\mu_2$,} \deleted{et le stabilis.} en esp. hom. \inserted{\unsure{sur|sous|$\to$} $\underline{\mathrm{SO}}(q)$, le stabilisateur d'une section est un ss-groupe fermé lisse à fibres conn.}%
% slanted

La décomposition en produit est vue avec l'équat. du cor. par ailleurs, notons seulement \deleted{\unsure{qu'on}} ,

\uline{Cor} La projection $\underline{\mathrm O}(q)\longrightarrow\mu_2$ est induite par dét. : $\underline{\mathrm O}(q)\longrightarrow\mathrm G_m$.

\deleted{Pour vérifier la lissité et la connexité de $\underline{\mathrm{SO}}$, \inserted{on est ramené à b).} \uline{\unsure{Obs|$\mathrm O(q)$}} qu'on a un $F\in G_{\{n\}}(E,q)$.}
% archive sheet number: 144
